
% BEGIN S2
 
\par\medskip\noindent\textbf{2 Preliminaries}\par
 
% BEGIN S2.p1
 
The syntax and inference rules of higher-order logic are recalled in Section 2.1. The basic notions about monad operators are summarized in Section 2.2.\par
 
% END S2.p1
 
% BEGIN S2.SS1
 
\par\medskip\noindent\textbf{2.1 Higher-Order Logic}\par
 
% BEGIN S2.SS1.p1
 
Higher-order logic is modeled using Church’s simple type theory [4]. Types are defined inductively: $\iota$ is the type of individuals, $o$ is the type of propositions, and if $\tau$ and $\sigma$ are types then $\tau\rightarrow\sigma$ is a type. For every type $\tau$, let $\mathcal{V}_{\tau}$ be the set of variables of type $\tau$ and $\mathcal{C}_{\tau}$ be a set of constants of type $\tau$. The set of variables $\mathcal{V}\mathrel{\mathop{:}}=\bigcup_{\tau}\mathcal{V}_{\tau}$ and the set of constants $\mathcal{C}\mathrel{\mathop{:}}=\bigcup_{\tau}\mathcal{C}_{\tau}$ are assumed to be disjoint. For any set of constants $\mathcal{C}$, the sets $\Lambda^{\mathcal{C}}_{\tau}$ of terms of type $\tau$ are defined by induction:\par
  
\par\noindent  $\blacksquare$ 
% BEGIN S2.I1.i1.p1
 
For every $x\in\mathcal{V}_{\tau}$, $x\in\Lambda^{\mathcal{C}}_{\tau}$.\par
 
% END S2.I1.i1.p1
 
 
\par\noindent  $\blacksquare$ 
% BEGIN S2.I1.i2.p1
 
For every $c\in\mathcal{C}_{\tau}$, $c\in\Lambda^{\mathcal{C}}_{\tau}$.\par
 
% END S2.I1.i2.p1
 
 
\par\noindent  $\blacksquare$ 
% BEGIN S2.I1.i3.p1
 
For every $x\in\mathcal{V}_{\tau}$ and $t\in\Lambda^{\mathcal{C}}_{\sigma}$, then $(\lambda x.\leavevmode\nobreak\ t)\in\Lambda^{\mathcal{C}}_{\tau\rightarrow\sigma}$.\par
 
% END S2.I1.i3.p1
 
 
\par\noindent  $\blacksquare$ 
% BEGIN S2.I1.i4.p1
 
For every $t\in\Lambda^{\mathcal{C}}_{\tau\rightarrow\sigma}$ and $u\in\Lambda^{\mathcal{C}}_{\tau}$, then $(t\leavevmode\nobreak\ u)\in\Lambda^{\mathcal{C}}_{\sigma}$.\par
 
% END S2.I1.i4.p1
 
  
$\lambda x.\leavevmode\nobreak\ t$ is a $\lambda$-abstraction and $t\leavevmode\nobreak\ u$ is an application. Computation is introduced in this $\lambda$-calculus thanks to the $\beta$-reduction rule $(\lambda x.\leavevmode\nobreak\ t)\leavevmode\nobreak\ u\hookrightarrow t[x\leftarrow u]$, where $t[x\leftarrow u]$ corresponds to the term $t$ in which $x$ has been substituted by $u$. We denote $\equiv_{\beta}$ the congruence generated by $\beta$-reduction. We write $\text{{FV}}(t_{1},\ldots,t_{n})$ for the set of free variables that occur in the terms $t_{1},\ldots,t_{n}$.\par
 
% END S2.SS1.p1
 
% BEGIN S2.SS1.p2
 
Formulas are terms of type $o$. There are particular constants defining the logical connectives and quantifiers: contradiction $\bot$ of type $o$, implication $\mathbin{\Rightarrow}$, conjunction $\mathbin{\wedge}$ and disjunction $\mathbin{\vee}$ of type $o\rightarrow o\rightarrow o$, and quantifiers ${\forall}_{\tau}$ and ${\exists}_{\tau}$ of type $(\tau\rightarrow o)\rightarrow o$. For convenience, terms of the form ${\forall}_{\tau}\leavevmode\nobreak\ (\lambda x.\leavevmode\nobreak\ A)$ and ${\exists}_{\tau}\leavevmode\nobreak\ (\lambda x.\leavevmode\nobreak\ A)$ are simply abbreviated as ${\forall}x.A$ and ${\exists}x.A$. The negation is defined by $\neg A\mathrel{\mathop{:}}=A\mathbin{\Rightarrow}\bot$ and the logical biconditional $\Leftrightarrow$ is defined by $A\Leftrightarrow B\mathrel{\mathop{:}}=(A\mathbin{\Rightarrow}B)\mathbin{\wedge}(B\mathbin{\Rightarrow}A)$. Contexts $\Gamma$ are finite sequences of formulas.\par
 
% END S2.SS1.p2
 
% BEGIN S2.SS1.p3
 
In the rest of this paper, we consider a logic $\mathsf{L}$ which is either minimal logic ($\mathsf{ML}$), intuitionistic logic ($\mathsf{IL}$), or classical logic ($\mathsf{CL}$). The natural deduction rules for $\mathsf{CL}$ are given in Figure 1. In $\mathsf{IL}$, we do not consider the principle of excluded middle PEM, and in $\mathsf{ML}$ we neither consider PEM nor the principle of explosion Bot-E. We write $\Gamma\vdash_{\mathsf{L}}A$ when $\Gamma\vdash A$ is derivable in the logic $\mathsf{L}$.\par
 
% END S2.SS1.p3
 
% BEGIN S2.F1

\begin{center}\includegraphics[width=0.95\linewidth,height=0.45\textheight,keepaspectratio]{prior-graphic-01.pdf}\end{center}
 
Figure 1: Natural deduction rules for classical logic.\par
 
% END S2.F1
 
% BEGIN S2.SS1.p4
 
We can also define, for every type $\tau$, an equality symbol $=_{\tau}$ of type $\tau\rightarrow\tau\rightarrow o$. The symbols are infix, and we write $t=u$ when there is no ambiguity on the type $\tau$. The natural deduction rules for equality are given in Figure 2. We write $\Gamma\vdash_{\mathsf{L}}^{*}A$ with $*\in\{\mathfrak{e},\mathfrak{e}\mathfrak{p},\mathfrak{e}\mathfrak{f},\mathfrak{e}\mathfrak{f}\mathfrak{p}\}$ when $\Gamma\vdash_{\mathsf{L}}A$ is derivable with additional inference rules: with Eq-I and Eq-E when $\mathfrak{e}$ occurs in $*$, with PropExt when $\mathfrak{p}$ occurs in $*$, and with FunExt when $\mathfrak{f}$ occurs in $*$.\par
 
% END S2.SS1.p4
 
% BEGIN S2.F2

\begin{center}\includegraphics[width=0.95\linewidth,height=0.45\textheight,keepaspectratio]{prior-graphic-02.pdf}\end{center}
 
Figure 2: Natural deduction rules for equality.\par
 
% END S2.F2
 
% END S2.SS1
 
% BEGIN S2.SS2
 
\par\medskip\noindent\textbf{2.2 Monad Operators}\par
 
% BEGIN S2.SS2.p1
 
We say that $T$ is a unary connective when it is a closed term of type $o\rightarrow o$. A unary connective $T$ is a monad operator of $\mathsf{L}$ when the judgments (unit) and (bind) hold for any formulas $A$ and $B$.\par
  
\par  \quad  $\displaystyle\vdash_{\mathsf{L}}A\mathbin{\Rightarrow}TA$\quad  \quad  (unit)\quad  
 
\par  \quad  $\displaystyle\vdash_{\mathsf{L}}(A\mathbin{\Rightarrow}TB)\mathbin{\Rightarrow}(TA\mathbin{\Rightarrow}TB)$\quad  \quad  (bind)\quad  
  
We use the terminology “monad operator” to avoid any confusion with monads in the sense of category theory. Monad operators are called lax modalities in [1] and nuclei in [22]. This definition of monad operators is equivalent to the axiomatization of strong monads in [6]. Following the Curry-Howard correspondence, monad operators coincide with the notion of monad in programming languages [15], for instance in Haskell.\par
 
% END S2.SS2.p1
 
% BEGIN Thmtheorem1
 
\par\medskip\noindent\textbf{Example 1.}\par
 
% BEGIN Thmtheorem1.p1
 
There are many monad operators of minimal logic:\par
  
\par\noindent  $\blacksquare$ 
% BEGIN S2.I2.i1.p1
 
$TA\mathrel{\mathop{:}}=A\mathbin{\vee}\bot$,\par
 
% END S2.I2.i1.p1
 
 
\par\noindent  $\blacksquare$ 
% BEGIN S2.I2.i2.p1
 
$T_{R}A\mathrel{\mathop{:}}=(A\mathbin{\Rightarrow}R)\mathbin{\Rightarrow}R$, with parameter $R$, which corresponds to the continuation monad,\par
 
% END S2.I2.i2.p1
 
 
\par\noindent  $\blacksquare$ 
% BEGIN S2.I2.i3.p1
 
$TA\mathrel{\mathop{:}}=\neg\neg A$, which is a special case of the previous one with $R\mathrel{\mathop{:}}=\bot$,\par
 
% END S2.I2.i3.p1
 
 
\par\noindent  $\blacksquare$ 
% BEGIN S2.I2.i4.p1
 
$T_{R}A\mathrel{\mathop{:}}=(A\mathbin{\Rightarrow}R)\mathbin{\Rightarrow}A$, with parameter $R$, which corresponds to the Peirce monad [6].\par
 
% END S2.I2.i4.p1
 
  
% END Thmtheorem1.p1
 
% END Thmtheorem1
 
% BEGIN S2.SS2.p2
 
We recall some formulas about monad operators that will be used in the rest of the paper.\par
 
% END S2.SS2.p2
 
% BEGIN Thmtheorem2
 
\par\medskip\noindent\textbf{Proposition 2.}\par
 
% BEGIN Thmtheorem2.p1
 
Let $T$ be a monad operator of $\mathsf{L}$, and $A$ and $B$ be formulas.\par
  
\par\noindent  1. 
% BEGIN S2.I3.i1.p1
 
$\vdash_{\mathsf{L}}TTA\Leftrightarrow TA$\par
 
% END S2.I3.i1.p1
 
 
\par\noindent  2. 
% BEGIN S2.I3.i2.p1
 
$\vdash_{\mathsf{L}}(A\mathbin{\Rightarrow}B)\mathbin{\Rightarrow}(TA\mathbin{\Rightarrow}TB)$\par
 
% END S2.I3.i2.p1
 
 
\par\noindent  3. 
% BEGIN S2.I3.i3.p1
 
$\vdash_{\mathsf{L}}T(A\mathbin{\wedge}B)\Leftrightarrow(TA\mathbin{\wedge}TB)$\par
 
% END S2.I3.i3.p1
 
 
\par\noindent  4. 
% BEGIN S2.I3.i4.p1
 
$\vdash_{\mathsf{L}}T(A\mathbin{\Rightarrow}B)\mathbin{\Rightarrow}(TA\mathbin{\Rightarrow}TB)$\par
 
% END S2.I3.i4.p1
 
 
\par\noindent  5. 
% BEGIN S2.I3.i5.p1
 
$\vdash_{\mathsf{L}}T(A\mathbin{\Rightarrow}TB)\Leftrightarrow(TA\mathbin{\Rightarrow}TB)$\par
 
% END S2.I3.i5.p1
 
 
\par\noindent  6. 
% BEGIN S2.I3.i6.p1
 
$\vdash_{\mathsf{L}}T(A\mathbin{\vee}B)\Leftrightarrow T(TA\mathbin{\vee}TB)$\par
 
% END S2.I3.i6.p1
 
 
\par\noindent  7. 
% BEGIN S2.I3.i7.p1
 
$\vdash_{\mathsf{L}}T({\forall}x.TA)\Leftrightarrow{\forall}x.TA$\par
 
% END S2.I3.i7.p1
 
 
\par\noindent  8. 
% BEGIN S2.I3.i8.p1
 
$\vdash_{\mathsf{L}}T({\exists}x.TA)\Leftrightarrow T({\exists}x.A)$\par
 
% END S2.I3.i8.p1
 
  
% END Thmtheorem2.p1
 
% END Thmtheorem2
 
% BEGIN S2.SS2.SSS0.P0.SPx1
 
\par\medskip\noindent\textbf{Proof.}\par
 
% BEGIN S2.SS2.SSS0.P0.SPx1.p1
 
We only show the most interesting cases.\par
  
\par\noindent  $\blacksquare$ 
% BEGIN S2.I4.i1.p1
 
Item 3 ($\mathbin{\Rightarrow}$): We have $\vdash_{\mathsf{L}}(A\mathbin{\wedge}B)\mathbin{\Rightarrow}A$ and $\vdash_{\mathsf{L}}(A\mathbin{\wedge}B)\mathbin{\Rightarrow}B$. Using Item 2 we easily derive $\vdash_{\mathsf{L}}T(A\mathbin{\wedge}B)\mathbin{\Rightarrow}(TA\mathbin{\wedge}TB)$.\par
 
% END S2.I4.i1.p1
 
 
\par\noindent  $\blacksquare$ 
% BEGIN S2.I4.i2.p1
 
Item 3 ($\Leftarrow$): Using (bind), we directly have $\vdash_{\mathsf{L}}(A\mathbin{\Rightarrow}T(A\mathbin{\wedge}B))\mathbin{\Rightarrow}TA\mathbin{\Rightarrow}T(A\mathbin{\wedge}B)$ and $\vdash_{\mathsf{L}}(B\mathbin{\Rightarrow}T(A\mathbin{\wedge}B))\mathbin{\Rightarrow}TB\mathbin{\Rightarrow}T(A\mathbin{\wedge}B)$. Using these two facts and (unit), we easily derive $\vdash_{\mathsf{L}}TA\mathbin{\Rightarrow}TB\mathbin{\Rightarrow}T(A\mathbin{\wedge}B)$. We conclude $\vdash_{\mathsf{L}}(TA\mathbin{\wedge}TB)\mathbin{\Rightarrow}T(A\mathbin{\wedge}B)$.\par
 
% END S2.I4.i2.p1
 
 
\par\noindent  $\blacksquare$ 
% BEGIN S2.I4.i3.p1
 
Item 4: We know that $\vdash_{\mathsf{L}}((A\mathbin{\Rightarrow}B)\mathbin{\wedge}A)\mathbin{\Rightarrow}B$. We derive $\vdash_{\mathsf{L}}T((A\mathbin{\Rightarrow}B)\mathbin{\wedge}A)\mathbin{\Rightarrow}TB$ using Item 2. By Item 3, we get $\vdash_{\mathsf{L}}(T(A\mathbin{\Rightarrow}B)\mathbin{\wedge}TA)\mathbin{\Rightarrow}TB$, therefore we conclude $\vdash_{\mathsf{L}}T(A\mathbin{\Rightarrow}B)\mathbin{\Rightarrow}(TA\mathbin{\Rightarrow}TB)$.\par
 
% END S2.I4.i3.p1
 
 
\par\noindent  $\blacksquare$ 
% BEGIN S2.I4.i4.p1
 
Item 5 ($\mathbin{\Rightarrow}$): We have $\vdash_{\mathsf{L}}T(A\mathbin{\Rightarrow}TB)\mathbin{\Rightarrow}(TA\mathbin{\Rightarrow}TTB)$ using Item 4. We conclude $\vdash_{\mathsf{L}}T(A\mathbin{\Rightarrow}TB)\mathbin{\Rightarrow}(TA\mathbin{\Rightarrow}TB)$ using Item 1.\par
 
% END S2.I4.i4.p1
 
 
\par\noindent  $\blacksquare$ 
% BEGIN S2.I4.i5.p1
 
Item 5 ($\Leftarrow$): Using (unit), we derive $\vdash_{\mathsf{L}}(TA\mathbin{\Rightarrow}TB)\mathbin{\Rightarrow}(A\mathbin{\Rightarrow}TB)$ and $\vdash_{\mathsf{L}}(A\mathbin{\Rightarrow}TB)\mathbin{\Rightarrow}T(A\mathbin{\Rightarrow}TB)$. It follows $\vdash_{\mathsf{L}}(TA\mathbin{\Rightarrow}TB)\mathbin{\Rightarrow}T(A\mathbin{\Rightarrow}TB)$.\par
 
% END S2.I4.i5.p1
 
 
\par\noindent  $\blacksquare$ 
% BEGIN S2.I4.i6.p1
 
Item 6 ($\Leftarrow$): Using (bind), we get $\vdash_{\mathsf{L}}(A\mathbin{\Rightarrow}T(A\mathbin{\vee}B))\mathbin{\Rightarrow}(TA\mathbin{\Rightarrow}T(A\mathbin{\vee}B))$ and $\vdash_{\mathsf{L}}(B\mathbin{\Rightarrow}T(A\mathbin{\vee}B))\mathbin{\Rightarrow}(TB\mathbin{\Rightarrow}T(A\mathbin{\vee}B))$. Using (unit), we derive $\vdash_{\mathsf{L}}TA\mathbin{\Rightarrow}T(A\mathbin{\vee}B)$ and $\vdash_{\mathsf{L}}TB\mathbin{\Rightarrow}T(A\mathbin{\vee}B)$. It follows $\vdash_{\mathsf{L}}(TA\mathbin{\vee}TB)\mathbin{\Rightarrow}T(A\mathbin{\vee}B)$. We conclude using (bind).\par
 
% END S2.I4.i6.p1
 
 
\par\noindent  $\blacksquare$ 
% BEGIN S2.I4.i7.p1
 
Item 7 ($\mathbin{\Rightarrow}$): We easily show $\vdash_{\mathsf{L}}({\forall}x.TA)\mathbin{\Rightarrow}TA[x\leftarrow y]$ for a fresh variable $y$. Using (bind) we get $\vdash_{\mathsf{L}}T({\forall}x.TA)\mathbin{\Rightarrow}TA[x\leftarrow y]$. We derive $T({\forall}x.TA)\vdash_{\mathsf{L}}TA[x\leftarrow y]$ and therefore $T({\forall}x.TA)\vdash_{\mathsf{L}}{\forall}x.TA$. We conclude $\vdash_{\mathsf{L}}T({\forall}x.TA)\mathbin{\Rightarrow}{\forall}x.TA$.\par
 
% END S2.I4.i7.p1
 
  
$\hfill\blacktriangleleft$\par
 
% END S2.SS2.SSS0.P0.SPx1.p1
 
% END S2.SS2.SSS0.P0.SPx1
 
% END S2.SS2
 
% END S2


% BEGIN S3
 
\par\medskip\noindent\textbf{3 Monad Translation}\par
 
% BEGIN S3.p1
 
The standard double-negation translations [13, 12, 10, 14] have been defined for first-order logic. Several translations generalize them to monad operators [1, 6, 22] and extend them to higher-order logic [3, 21]. We want to define a translation that does both.\par
 
% END S3.p1
 
% BEGIN S3.p2
 
Kolmogorov’s translation [13] and the Gödel-Gentzen translation [12, 10] cannot be extended to higher-order logic directly [3], because they do not preserve $\beta$-conversion. For example, let us take the Gödel-Gentzen translation of $(\lambda R.\leavevmode\nobreak\ R\mathbin{\wedge}Q)\leavevmode\nobreak\ P\hookrightarrow P\mathbin{\wedge}Q$. The translation of $(\lambda R.\leavevmode\nobreak\ R\mathbin{\wedge}Q)\leavevmode\nobreak\ P$ is $(\lambda R.\leavevmode\nobreak\ \neg\neg R\mathbin{\wedge}\neg\neg Q)\leavevmode\nobreak\ \neg\neg P$, which $\beta$-reduces to $\neg\neg\neg\neg P\mathbin{\wedge}\neg\neg Q$, while the translation of $P\mathbin{\wedge}Q$ is $\neg\neg P\mathbin{\wedge}\neg\neg Q$. Unlike Kolmogorov’s translation and the Gödel-Gentzen translation, Kuroda’s translation [14] preserves $\beta$-conversion and can be directly extended to higher-order logic [3, 21].\par
 
% END S3.p2
 
% BEGIN S3.p3
 
Moreover, a Kuroda-style monad translation [22] has been defined for first-order logic. It corresponds to Kuroda’s translation, in which we insert monad operators instead of double negations. It also introduces additional monad operators after implications, and such modification is necessary to generalize Kuroda’s translation to any monad operator. We will see in Section 7 how to remove this constraint for particular monad operators.\par
 
% END S3.p3
 
% BEGIN S3.p4
 
We define here a Kuroda-style monad translation for higher-order logic, taking advantage of the ideas developed in [21] and [22].\par
 
% END S3.p4
 
% BEGIN Thmtheorem3
 
\par\medskip\noindent\textbf{Definition 3 (Monad translation for higher-order logic).}\par
 
% BEGIN Thmtheorem3.p1
 
Let $A$ be a formula in higher-order logic and $T$ be a monad operator. Its monad translation is $A^{T}\mathrel{\mathop{:}}=TA_{T}$, where $t\mapsto t_{T}$ is inductively defined by:\par
  
\par  \quad  \[\begin{array}[]{lll}x_{T}&\mathrel{\mathop{:}}=&x\\
c_{T}&\mathrel{\mathop{:}}=&\left\{\begin{array}[]{ll}\lambda p.\leavevmode\nobreak\ {\forall}x.T(p\leavevmode\nobreak\ x)&\text{ if $c={\forall}$}\\
\lambda p.\leavevmode\nobreak\ \lambda q.\leavevmode\nobreak\ p\mathbin{\Rightarrow}Tq&\text{ if $c=\mathbin{\Rightarrow}$}\\
c&\text{ otherwise}\\
\end{array}\right.\\
(\lambda x.\leavevmode\nobreak\ t)_{T}&\mathrel{\mathop{:}}=&\lambda x.\leavevmode\nobreak\ t_{T}\\
(t\leavevmode\nobreak\ u)_{T}&\mathrel{\mathop{:}}=&t_{T}\leavevmode\nobreak\ u_{T}\\
\end{array}\]\quad  \quad  
  
% END Thmtheorem3.p1
 
% END Thmtheorem3
 
% BEGIN S3.p5
 
In first-order logic, we have $(A[z\leftarrow w])^{T}=A^{T}[z\leftarrow w]$. As we are in higher-order logic, the substituted term $w$ may be modified by the translation. We would intuitively expect $(A[z\leftarrow w])^{T}=A^{T}[z\leftarrow w^{T}]$, but we actually have $(A[z\leftarrow w])^{T}=A^{T}[z\leftarrow w_{T}]$, because the monad operator in front of the formula is inserted after the inductive translation. It follows that the translation $t\mapsto t_{T}$ commutes with substitution, but not $A\mapsto A^{T}$. Both translations $t\mapsto t_{T}$ and $A\mapsto A^{T}$ preserve convertibility.\par
 
% END S3.p5
 
% BEGIN Thmtheorem4
 
\par\medskip\noindent\textbf{Proposition 4.}\par
 
% BEGIN Thmtheorem4.p1
 
For any term $t$, we have $(t[z\leftarrow w])_{T}=t_{T}[z\leftarrow w_{T}]$.\par
 
% END Thmtheorem4.p1
 
% END Thmtheorem4
 
% BEGIN S3.SS0.SSS0.P0.SPx1
 
\par\medskip\noindent\textbf{Proof.}\par
 
% BEGIN S3.SS0.SSS0.P0.SPx1.p1
 
By induction on the term $t$. $\hfill\blacktriangleleft$\par
 
% END S3.SS0.SSS0.P0.SPx1.p1
 
% BEGIN Thmtheorem5
 
\par\medskip\noindent\textbf{Corollary 5.}\par
 
% BEGIN Thmtheorem5.p1
 
For any higher-order formula $A$, we have $(A[z\leftarrow w])^{T}=A^{T}[z\leftarrow w_{T}]$.\par
 
% END Thmtheorem5.p1
 
% END Thmtheorem5
 
% BEGIN S3.SS0.SSS0.P0.SPx1.p2
 
In higher-order logic, terms are considered modulo $\beta$-convertibility. It follows that the translation must preserve $\beta$-convertibility for it to preserve provability.\par
 
% END S3.SS0.SSS0.P0.SPx1.p2
 
% BEGIN Thmtheorem6
 
\par\medskip\noindent\textbf{Proposition 6.}\par
 
% BEGIN Thmtheorem6.p1
 
For any terms $t$ and $u$, if $t\equiv_{\beta}u$ then $t_{T}\equiv_{\beta}u_{T}$.\par
 
% END Thmtheorem6.p1
 
% END Thmtheorem6
 
% END S3.SS0.SSS0.P0.SPx1
 
% BEGIN S3.SS0.SSS0.P0.SPx2
 
\par\medskip\noindent\textbf{Proof.}\par
 
% BEGIN S3.SS0.SSS0.P0.SPx2.p1
 
By definition $((\lambda x.\leavevmode\nobreak\ t)\leavevmode\nobreak\ u)_{T}=(\lambda x.\leavevmode\nobreak\ t_{T})\leavevmode\nobreak\ u_{T}$, and we have $(\lambda x.\leavevmode\nobreak\ t_{T})\leavevmode\nobreak\ u_{T}\hookrightarrow t_{T}[x\leftarrow u_{T}]$. Using Proposition 4, we get $((\lambda x.\leavevmode\nobreak\ t)\leavevmode\nobreak\ u)_{T}\equiv_{\beta}(t[x\leftarrow u])_{T}$. Closure by context, reflexivity, symmetry, and transitivity are immediate. $\hfill\blacktriangleleft$\par
 
% END S3.SS0.SSS0.P0.SPx2.p1
 
% BEGIN Thmtheorem7
 
\par\medskip\noindent\textbf{Corollary 7.}\par
 
% BEGIN Thmtheorem7.p1
 
For any higher-order formulas $A$ and $B$, if $A\equiv_{\beta}B$ then $A^{T}\equiv_{\beta}B^{T}$.\par
 
% END Thmtheorem7.p1
 
% END Thmtheorem7
 
% END S3.SS0.SSS0.P0.SPx2
 
% END S3


% BEGIN S4
 
\par\medskip\noindent\textbf{4 Monad Embedding}\par
 
% BEGIN S4.p1
 
The monad translation generalizes double-negation translations with monad operators instead of double negations. Likewise, the monad translation eliminates the use of a particular formula [1, 6], called $T$-elimination, while double-negation translations eliminate the use of the principle of excluded middle.\par
  
\par  \quad  $\displaystyle TA\mathbin{\Rightarrow}A$\quad  \quad  ($T$-elimination)\quad  
  
More formally, we extend the logic $\mathsf{L}$ with the inference rule T-Elim, and we write $\Gamma\vdash_{\mathsf{L}+\mathsf{T}}A$ when $\Gamma\vdash A$ is derivable using the rules of $\mathsf{L}$ and T-Elim.\par
 
% END S4.p1
 
% BEGIN S4.p2
 
\begin{center}\includegraphics[width=0.95\linewidth,height=0.45\textheight,keepaspectratio]{prior-graphic-03.pdf}\end{center}
 
% END S4.p2
 
% BEGIN S4.p3
 
The monad translation eliminates any use of the T-Elim inference rule–this is the soundness property. The characterization property states that any formula and its translation are equivalent assuming T-Elim.\par
 
% END S4.p3
 
% BEGIN S4.p4
 
When extending the double-negation elimination to higher-order logic [21], the characterization property does not trivially hold, but does hold when we consider functional extensionality and propositional extensionality. The same reasoning applies for the extension of the monad translation to higher-order logic.\par
 
% END S4.p4
 
% BEGIN Thmtheorem8
 
\par\medskip\noindent\textbf{Theorem 8 (Monad embedding).}\par
 
% BEGIN Thmtheorem8.p1
 
Let $T$ be a monad operator of $\mathsf{L}$ such that we have $\vdash_{\mathsf{L}}(TA)_{T}\Leftrightarrow TA_{T}$ for any formula $A$. Let $\Gamma,A$ be higher-order formulas.\par
  
\par\noindent  1. 
% BEGIN S4.I1.i1.p1
 
If $\Gamma\vdash_{\mathsf{L}+\mathsf{T}}A$ then $\Gamma_{T}\vdash_{\mathsf{L}}A^{T}$.\par
 
% END S4.I1.i1.p1
 
 
\par\noindent  2. 
% BEGIN S4.I1.i2.p1
 
$\vdash_{\mathsf{L}+\mathsf{T}}^{\mathfrak{e}\mathfrak{f}\mathfrak{p}}A^{T}\Leftrightarrow
A$.\par
 
% END S4.I1.i2.p1
 
  
% END Thmtheorem8.p1
 
% END Thmtheorem8
 
% BEGIN Thmtheorem9
 
\par\medskip\noindent\textbf{$\blacktriangleright$ Remark 9.}\par
 
% BEGIN Thmtheorem9.p1
 
For Kolmogorov’s translation and the Gödel-Gentzen translation, the soundness property is stated with $\Gamma^{T}$ instead of $\Gamma_{T}$. When considering Kuroda’s translation, we can actually omit the monad operators at the head of the formulas of the context. The version with $\Gamma^{T}$ is provable as well.\par
 
% END Thmtheorem9.p1
 
% END Thmtheorem9
 
% BEGIN S4.SS0.SSS0.P0.SPx1
 
\par\medskip\noindent\textbf{Proof of Theorem 8(1).}\par
 
% BEGIN S4.SS0.SSS0.P0.SPx1.p1
 
We proceed by induction on the derivation. Most of the cases are direct applications of Proposition 2. Conv derives from Corollary 7. We only show the most interesting cases:\par
  
\par\noindent  $\blacksquare$ 
% BEGIN S4.I2.i1.p1
 
Or-E: Suppose that $\Gamma_{T}\vdash_{\mathsf{L}}T(A_{T}\mathbin{\vee}B_{T})$ and $\Gamma_{T},A_{T}\vdash_{\mathsf{L}}TC_{T}$ and $\Gamma_{T},B_{T}\vdash_{\mathsf{L}}TC_{T}$. To prove $\Gamma_{T}\vdash_{\mathsf{L}}TC_{T}$, it suffices to have $\Gamma_{T}\vdash_{\mathsf{L}}T(A_{T}\mathbin{\vee}B_{T})\mathbin{\Rightarrow}TC_{T}$.\par
 
% END S4.I2.i1.p1
 
% BEGIN S4.I2.i1.p2
 
Using the second and third hypotheses, we derive $\Gamma_{T}\vdash_{\mathsf{L}}(A_{T}\mathbin{\vee}B_{T})\mathbin{\Rightarrow}TC_{T}$. Using (bind) we get $\Gamma_{T}\vdash_{\mathsf{L}}T(A_{T}\mathbin{\vee}B_{T})\mathbin{\Rightarrow}TC_{T}$ and we conclude the case.\par
 
% END S4.I2.i1.p2
 
 
\par\noindent  $\blacksquare$ 
% BEGIN S4.I2.i2.p1
 
All-I: Suppose $\Gamma_{T}\vdash_{\mathsf{L}}TA_{T}$. We derive $\Gamma_{T}\vdash_{\mathsf{L}}{\forall}x.TA^{T}$, that is $\Gamma_{T}\vdash_{\mathsf{L}}({\forall}x.A)_{T}$. We conclude using (unit).\par
 
% END S4.I2.i2.p1
 
 
\par\noindent  $\blacksquare$ 
% BEGIN S4.I2.i3.p1
 
All-E: Suppose $\Gamma_{T}\vdash_{\mathsf{L}}T({\forall}x.TA_{T})$. Using Proposition 2(7), we get $\Gamma_{T}\vdash_{\mathsf{L}}{\forall}x.TA_{T}$. From All-E we obtain $\Gamma_{T}\vdash_{\mathsf{L}}(TA_{T})[x\leftarrow t_{T}]$. We conclude with Corollary 5.\par
 
% END S4.I2.i3.p1
 
 
\par\noindent  $\blacksquare$ 
% BEGIN S4.I2.i4.p1
 
Ex-I: Suppose $\Gamma_{T}\vdash_{\mathsf{L}}T(A[x\leftarrow t])_{T}$. Using Corollary 5, we get $\Gamma_{T}\vdash_{\mathsf{L}}TA_{T}[x\leftarrow t_{T}]$. From Ex-I we obtain $\Gamma_{T}\vdash_{\mathsf{L}}{\exists}x.TA_{T}$. We conclude using (unit) and Proposition 2(8).\par
 
% END S4.I2.i4.p1
 
 
\par\noindent  $\blacksquare$ 
% BEGIN S4.I2.i5.p1
 
Ex-E: Suppose $\Gamma_{T}\vdash_{\mathsf{L}}T({\exists}x.A_{T})$ and $\Gamma_{T},A_{T}\vdash_{\mathsf{L}}TC_{T}$. To conclude the proof, it suffices to have $\Gamma_{T}\vdash_{\mathsf{L}}T({\exists}x.A_{T})\mathbin{\Rightarrow}TC_{T}$.\par
 
% END S4.I2.i5.p1
 
% BEGIN S4.I2.i5.p2
 
We have $\Gamma_{T},{\exists}x.A_{T}\vdash_{\mathsf{L}}{\exists}x.A_{T}$ (by Ax) and $\Gamma_{T},{\exists}x.A_{T},A_{T}\vdash_{\mathsf{L}}TC_{T}$ (by weakening of the second hypothesis). Using Ex-E, we derive $\Gamma_{T},{\exists}x.A_{T}\vdash_{\mathsf{L}}TC_{T}$ and therefore $\Gamma_{T}\vdash_{\mathsf{L}}({\exists}x.A_{T})\mathbin{\Rightarrow}TC_{T}$. Using (bind), we obtain $\Gamma_{T}\vdash_{\mathsf{L}}T({\exists}x.A_{T})\mathbin{\Rightarrow}TC_{T}$ and we conclude the case.\par
 
% END S4.I2.i5.p2
 
 
\par\noindent  $\blacksquare$ 
% BEGIN S4.I2.i6.p1
 
Bot-E: Suppose $\Gamma_{T}\vdash_{\mathsf{IL}}T\bot$. As we are in intuitionistic logic, we have $\Gamma_{T}\vdash_{\mathsf{IL}}\bot\mathbin{\Rightarrow}A_{T}$. We deduce $\Gamma_{T}\vdash_{\mathsf{IL}}T\bot\mathbin{\Rightarrow}TA_{T}$ using Proposition 2(2). We conclude $\Gamma_{T}\vdash_{\mathsf{IL}}TA_{T}$.\par
 
% END S4.I2.i6.p1
 
 
\par\noindent  $\blacksquare$ 
% BEGIN S4.I2.i7.p1
 
PEM: As we are in classical logic, we have $\Gamma_{T}\vdash_{\mathsf{CL}}A_{T}\mathbin{\vee}(A_{T}\mathbin{\Rightarrow}\bot)$. Using (unit), we derive $\Gamma_{T}\vdash_{\mathsf{CL}}A_{T}\mathbin{\vee}(A_{T}\mathbin{\Rightarrow}T\bot)$ and then we conclude $\Gamma_{T}\vdash_{\mathsf{CL}}T(A_{T}\mathbin{\vee}(A_{T}\mathbin{\Rightarrow}T\bot))$.\par
 
% END S4.I2.i7.p1
 
 
\par\noindent  $\blacksquare$ 
% BEGIN S4.I2.i8.p1
 
T-Elim: We have $(TA\mathbin{\Rightarrow}A)_{T}=((TA)_{T}\mathbin{\Rightarrow}TA_{T})$. We know $\vdash_{\mathsf{L}}(TA)_{T}\Leftrightarrow TA_{T}$, so we derive $\Gamma_{T}\vdash_{\mathsf{L}}(TA\mathbin{\Rightarrow}A)_{T}$. We conclude $\Gamma_{T}\vdash_{\mathsf{L}}T(TA\mathbin{\Rightarrow}A)_{T}$ using (unit).\par
 
% END S4.I2.i8.p1
 
  
$\hfill\blacktriangleleft$\par
 
% END S4.SS0.SSS0.P0.SPx1.p1
 
% BEGIN S4.SS0.SSS0.P0.SPx1.p2
 
Before proving the second item of Theorem 8, we show the following lemma, stating that any term and its translation are equal, assuming T-Elim, functional extensionality and propositional extensionality.\par
 
% END S4.SS0.SSS0.P0.SPx1.p2
 
% BEGIN Thmtheorem10
 
\par\medskip\noindent\textbf{Lemma 10.}\par
 
% BEGIN Thmtheorem10.p1
 
For any term $t$, we have $\vdash_{\mathsf{L}+\mathsf{T}}^{\mathfrak{e}\mathfrak{f}\mathfrak{p}}t_{T}=t$.\par
 
% END Thmtheorem10.p1
 
% END Thmtheorem10
 
% END S4.SS0.SSS0.P0.SPx1
 
% BEGIN S4.SS0.SSS0.P0.SPx2
 
\par\medskip\noindent\textbf{Proof.}\par
 
% BEGIN S4.SS0.SSS0.P0.SPx2.p1
 
We proceed by induction on the term $t$. We directly have $\vdash_{\mathsf{L}+\mathsf{T}}^{\mathfrak{e}\mathfrak{f}\mathfrak{p}}x_{T}=x$ and $\vdash_{\mathsf{L}+\mathsf{T}}^{\mathfrak{e}\mathfrak{f}\mathfrak{p}}c_{T}=c$ for $c\notin\{{\forall},\mathbin{\Rightarrow}\}$. For $c\in\{{\forall},\mathbin{\Rightarrow}\}$, we derive $\vdash_{\mathsf{L}+\mathsf{T}}^{\mathfrak{e}\mathfrak{f}\mathfrak{p}}c_{T}=c$ from PropExt, FunExt, T-Elim and (unit). We have $\vdash_{\mathsf{L}+\mathsf{T}}^{\mathfrak{e}\mathfrak{f}\mathfrak{p}}(t\leavevmode\nobreak\ u)_{T}=t\leavevmode\nobreak\ u$ using the induction hypotheses and Eq-E. We derive $\vdash_{\mathsf{L}+\mathsf{T}}^{\mathfrak{e}\mathfrak{f}\mathfrak{p}}(\lambda x.\leavevmode\nobreak\ t)_{T}=\lambda x.\leavevmode\nobreak\ t$ using the induction hypothesis and FunExt. $\hfill\blacktriangleleft$\par
 
% END S4.SS0.SSS0.P0.SPx2.p1
 
% END S4.SS0.SSS0.P0.SPx2
 
% BEGIN S4.SS0.SSS0.P0.SPx3
 
\par\medskip\noindent\textbf{Proof of Theorem 8(2).}\par
 
% BEGIN S4.SS0.SSS0.P0.SPx3.p1
 
We derive $\vdash_{\mathsf{L}+\mathsf{T}}^{\mathfrak{e}\mathfrak{f}\mathfrak{p}}A_{T}\Leftrightarrow
A$ using Lemma 10, and then we conclude $\vdash_{\mathsf{L}+\mathsf{T}}^{\mathfrak{e}\mathfrak{f}\mathfrak{p}}A^{T}\Leftrightarrow
A$ using (unit) and T-Elim. $\hfill\blacktriangleleft$\par
 
% END S4.SS0.SSS0.P0.SPx3.p1
 
% BEGIN S4.SS0.SSS0.P0.SPx3.p2
 
We proved that the monad translation satisfies the soundness and characterization properties. For the characterization property, we had to introduce symbols of equality and to assume both functional extensionality and propositional extensionality. But does the soundness property still hold in the presence of equality?\par
 
% END S4.SS0.SSS0.P0.SPx3.p2
 
% BEGIN S4.SS0.SSS0.P0.SPx3.p3
 
The proof of the soundness property does not naturally extend when we consider the symbols and inference rules of equality. Let us consider the formulas\par
  
\par  \quad  $\displaystyle{\forall}x{\forall}y.\leavevmode\nobreak\ T(x=y)\mathbin{\Rightarrow}x=y$\quad  \quad  ($\Delta^{\mathfrak{e}\mathfrak{f}}$)\quad  
 
\par  \quad  $\displaystyle{\forall}x{\forall}y.\leavevmode\nobreak\ (Tx=Ty)\mathbin{\Rightarrow}x=y$\quad  \quad  ($\Delta^{\mathfrak{e}\mathfrak{p}}$)\quad  
  
regarding $T$ and equality predicates. We write $\Delta^{\mathfrak{e}\mathfrak{f}\mathfrak{p}}$ for the context $\Delta^{\mathfrak{e}\mathfrak{f}},\Delta^{\mathfrak{e}\mathfrak{p}}$. When proving the soundness property in the presence of equality, the first formula $\Delta^{\mathfrak{e}\mathfrak{f}}$ is useful for the FunExt case and the second formula $\Delta^{\mathfrak{e}\mathfrak{p}}$ is useful for the PropExt case.\par
 
% END S4.SS0.SSS0.P0.SPx3.p3
 
% BEGIN Thmtheorem11
 
\par\medskip\noindent\textbf{Theorem 11 (Soundness for equality).}\par
 
% BEGIN Thmtheorem11.p1
 
Let $T$ be a monad operator of $\mathsf{L}$ such that we have $\vdash_{\mathsf{L}}(TA)_{T}\Leftrightarrow TA_{T}$ for any formula $A$. Let $\Gamma,A$ be higher-order formulas.\par
  
\par\noindent  1. 
% BEGIN S4.I3.i1.p1
 
If $\Gamma\vdash_{\mathsf{L}+\mathsf{T}}^{\mathfrak{e}}A$ then $\Gamma_{T}\vdash_{\mathsf{L}}^{\mathfrak{e}}A^{T}$.\par
 
% END S4.I3.i1.p1
 
 
\par\noindent  2. 
% BEGIN S4.I3.i2.p1
 
For $*\in\{\mathfrak{e}\mathfrak{f},\mathfrak{e}\mathfrak{p},\mathfrak{e}\mathfrak{f}\mathfrak{p}\}$, if $\Gamma\vdash_{\mathsf{L}+\mathsf{T}}^{*}A$ then $\Delta^{*},\Gamma_{T}\vdash_{\mathsf{L}}^{*}A^{T}$.\par
 
% END S4.I3.i2.p1
 
  
% END Thmtheorem11.p1
 
% END Thmtheorem11
 
% END S4.SS0.SSS0.P0.SPx3
 
% BEGIN S4.SS0.SSS0.P0.SPx4
 
\par\medskip\noindent\textbf{Proof.}\par
 
% BEGIN S4.SS0.SSS0.P0.SPx4.p1
 
For the first item, we complete the proof of Theorem 8 with the following cases:\par
  
\par\noindent  $\blacksquare$ 
% BEGIN S4.I4.i1.p1
 
Eq-I: We directly use Eq-I and (unit).\par
 
% END S4.I4.i1.p1
 
 
\par\noindent  $\blacksquare$ 
% BEGIN S4.I4.i2.p1
 
Eq-E: Suppose $\Gamma_{T}\vdash_{\mathsf{L}}^{*}TA_{T}[x\leftarrow u_{T}]$ and $\Gamma_{T}\vdash_{\mathsf{L}}^{*}T(u_{T}=v_{T})$. Using Eq-E, we derive $\vdash_{\mathsf{L}}^{*}A_{T}[x\leftarrow u_{T}]\mathbin{\Rightarrow}(u_{T}=v_{T})\mathbin{\Rightarrow}A_{T}[x\leftarrow v_{T}]$. Therefore, we get $\vdash_{\mathsf{L}}^{*}TA_{T}[x\leftarrow u_{T}]\mathbin{\Rightarrow}T(u_{T}=v_{T})\mathbin{\Rightarrow}TA_{T}[x\leftarrow v_{T}]$ using Proposition 2(2) and Proposition 2(4). We conclude $\Gamma_{T}\vdash_{\mathsf{L}}^{*}TA_{T}[x\leftarrow v_{T}]$.\par
 
% END S4.I4.i2.p1
 
  
For the second item, most cases are those of the first item that are reused using weakening on $\Delta^{*}$. The remaining cases are:\par
  
\par\noindent  $\blacksquare$ 
% BEGIN S4.I5.i1.p1
 
FunExt: Suppose $\Delta^{*},\Gamma_{T}\vdash_{\mathsf{L}}^{*}T(f_{T}\leavevmode\nobreak\ x=g_{T}\leavevmode\nobreak\ x)$. Since $x\notin\text{{FV}}(\Gamma,f,g)$, we directly have $x\notin\text{{FV}}(\Gamma_{T},f_{T},g_{T})$. We get $\Delta^{*},\Gamma_{T}\vdash_{\mathsf{L}}^{*}f_{T}\leavevmode\nobreak\ x=g_{T}\leavevmode\nobreak\ x$ using $\Delta^{\mathfrak{e}\mathfrak{f}}$. We conclude using FunExt and (unit).\par
 
% END S4.I5.i1.p1
 
 
\par\noindent  $\blacksquare$ 
% BEGIN S4.I5.i2.p1
 
PropExt: Suppose $\Delta^{*},\Gamma_{T}\vdash_{\mathsf{L}}^{*}T(A_{T}\mathbin{\Rightarrow}TB_{T})$ and $\Delta^{*},\Gamma_{T}\vdash_{\mathsf{L}}^{*}T(B_{T}\mathbin{\Rightarrow}TA_{T})$. Using Proposition 2(5), we get $\Delta^{*},\Gamma_{T}\vdash_{\mathsf{L}}^{*}TA_{T}\mathbin{\Rightarrow}TB_{T}$ and $\Delta^{*},\Gamma_{T}\vdash_{\mathsf{L}}^{*}TB_{T}\mathbin{\Rightarrow}TA_{T}$. Using PropExt we derive $\Delta^{*},\Gamma_{T}\vdash_{\mathsf{L}}^{*}TA_{T}=TB_{T}$. We conclude using $\Delta^{\mathfrak{e}\mathfrak{p}}$ and (unit).\par
 
% END S4.I5.i2.p1
 
  
$\hfill\blacktriangleleft$\par
 
% END S4.SS0.SSS0.P0.SPx4.p1
 
% BEGIN S4.SS0.SSS0.P0.SPx4.p2
 
When extending Kuroda’s double-negation translation to higher-order logic [21], we only considered the formula $\Delta^{\mathfrak{e}\mathfrak{f}}$. The generalization to any monad operator requires us to additionally consider the formula $\Delta^{\mathfrak{e}\mathfrak{p}}$. In Section 7, we will see that this additional assumption can be dropped when considering factorizable monad operators.\par
 
% END S4.SS0.SSS0.P0.SPx4.p2
 
% END S4.SS0.SSS0.P0.SPx4
 
% END S4


% BEGIN S5
 
\par\medskip\noindent\textbf{5 Embeddings between Logics}\par
 
% BEGIN S5.p1
 
Such monad translation can be directly applied with particular monad operators to embed classical logic into intuitionistic logic or intuitionistic logic into minimal logic.\par
 
% END S5.p1
 
% BEGIN S5.p2
 
There are many classical laws, such as the principle of excluded middle, the double-negation elimination, Peirce’s law, or Clavius’s law. All these laws are equivalent in intuitionistic logic.\par
  
\par  \quad  $\displaystyle\neg\neg A\mathbin{\Rightarrow}A$\quad  \quad  (double-negation elimination)\quad  
 
\par  \quad  $\displaystyle(\neg A\mathbin{\Rightarrow}A)\mathbin{\Rightarrow}A$\quad  \quad  (Clavius’s law)\quad  
 
\par  \quad  $\displaystyle((A\mathbin{\Rightarrow}B)\mathbin{\Rightarrow}A)\mathbin{\Rightarrow}A$\quad  \quad  (Peirce’s law)\quad  
  
The $T$-elimination principle corresponds to the double-negation elimination when we choose $TA\mathrel{\mathop{:}}=(A\mathbin{\Rightarrow}\bot)\mathbin{\Rightarrow}\bot$ and to Clavius’s law when we choose $TA\mathrel{\mathop{:}}=(A\mathbin{\Rightarrow}\bot)\mathbin{\Rightarrow}A$. For these monad operators, the T-Elim inference rule is therefore equivalent to PEM, and eliminating any use of T-elim in the derivations means transforming classical provability into intuitionistic provability. When $T_{R}A\mathrel{\mathop{:}}=(A\mathbin{\Rightarrow}R)\mathbin{\Rightarrow}A$, the $T$-elimination principle only corresponds to an instance of Peirce’s law with $B\mathrel{\mathop{:}}=R$.\par
 
% END S5.p2
 
% BEGIN Thmtheorem12
 
\par\medskip\noindent\textbf{Corollary 12 (Embeddings of classical logic into intuitionistic logic).}\par
 
% BEGIN Thmtheorem12.p1
 
Let $\Gamma,A$ be higher-order formulas.\par
  
\par\noindent  1. 
% BEGIN S5.I1.i1.p1
 
If $\Gamma\vdash_{\mathsf{CL}}A$ then $\Gamma_{T}\vdash_{\mathsf{IL}}A^{T}$ with $TA\mathrel{\mathop{:}}=(A\mathbin{\Rightarrow}\bot)\mathbin{\Rightarrow}\bot$.\par
 
% END S5.I1.i1.p1
 
 
\par\noindent  2. 
% BEGIN S5.I1.i2.p1
 
If $\Gamma\vdash_{\mathsf{CL}}A$ then $\Gamma_{T}\vdash_{\mathsf{IL}}A^{T}$ with $TA\mathrel{\mathop{:}}=(A\mathbin{\Rightarrow}\bot)\mathbin{\Rightarrow}A$.\par
 
% END S5.I1.i2.p1
 
 
\par\noindent  3. 
% BEGIN S5.I1.i3.p1
 
If $\Gamma\vdash_{\mathsf{CL}}A$ then there exists a formula $R$ such that $\Gamma_{T_{R}}\vdash_{\mathsf{IL}}A^{T_{R}}$ with $T_{R}A\mathrel{\mathop{:}}=(A\mathbin{\Rightarrow}R)\mathbin{\Rightarrow}A$.\par
 
% END S5.I1.i3.p1
 
  
Moreover, for each of the above, we have $\vdash_{\mathsf{CL}}^{\mathfrak{e}\mathfrak{f}\mathfrak{p}}A^{T}\Leftrightarrow
A$.\par
 
% END Thmtheorem12.p1
 
% END Thmtheorem12
 
% BEGIN S5.SS0.SSS0.P0.SPx1
 
\par\medskip\noindent\textbf{Proof.}\par
 
% BEGIN S5.SS0.SSS0.P0.SPx1.p1
 
We apply Theorem 8 with $\mathsf{L}\mathrel{\mathop{:}}=\mathsf{IL}$. T-Elim is equivalent to PEM, so $\mathsf{IL}+\mathsf{T}$ corresponds to $\mathsf{CL}$. For the third item, we have to give the correct parameter $R$. The derivation $\Gamma\vdash_{\mathsf{CL}}A$ is finite, so Peirce’s law is used a finite number of times with formulas $B_{1},\ldots,B_{n}$ of free variables $\vec{x_{1}},\ldots,\vec{x_{n}}$. Such free variables could be captured later in the derivation. As remarked by Escardó and Oliva [6], $T_{R_{1}\mathbin{\wedge}R_{2}}A\mathbin{\Rightarrow}A\vdash_{\mathsf{IL}}(T_{R_{1}}A\mathbin{\Rightarrow}A)\mathbin{\wedge}(T_{R_{2}}A\mathbin{\Rightarrow}A)$. With the additional remark that $T_{{\forall}x.R}A\mathbin{\Rightarrow}A\vdash_{\mathsf{IL}}T_{R[x\leftarrow t]}A\mathbin{\Rightarrow}A$, we define $R$ to be the conjunction ${\forall}\vec{x_{1}}.B_{1}\mathbin{\wedge}\ldots\mathbin{\wedge}{\forall}\vec{x_{n}}.B_{n}$, so that the monad translation eliminates the $n$ instances of Peirce’s law that occur in the derivation. $\hfill\blacktriangleleft$\par
 
% END S5.SS0.SSS0.P0.SPx1.p1
 
% BEGIN S5.SS0.SSS0.P0.SPx1.p2
 
When $TA\mathrel{\mathop{:}}=A\mathbin{\vee}\bot$, the $T$-elimination principle is equivalent to the principle of explosion, and the T-Elim inference rule is therefore equivalent to the Bot-E inference rule. In that case, eliminating any use of T-elim in the derivations means transforming intuitionistic provability into minimal provability.\par
 
% END S5.SS0.SSS0.P0.SPx1.p2
 
% BEGIN Thmtheorem13
 
\par\medskip\noindent\textbf{Corollary 13 (Embedding of intuitionistic logic into minimal logic).}\par
 
% BEGIN Thmtheorem13.p1
 
Let $\Gamma,A$ be higher-order formulas. If $\Gamma\vdash_{\mathsf{IL}}A$ then $\Gamma_{T}\vdash_{\mathsf{ML}}A^{T}$ with $TA\mathrel{\mathop{:}}=A\mathbin{\vee}\bot$. Moreover, we have $\vdash_{\mathsf{IL}}^{\mathfrak{e}\mathfrak{f}\mathfrak{p}}A^{T}\Leftrightarrow
A$.\par
 
% END Thmtheorem13.p1
 
% END Thmtheorem13
 
% END S5.SS0.SSS0.P0.SPx1
 
% BEGIN S5.SS0.SSS0.P0.SPx2
 
\par\medskip\noindent\textbf{Proof.}\par
 
% BEGIN S5.SS0.SSS0.P0.SPx2.p1
 
We apply Theorem 8 with $\mathsf{L}\mathrel{\mathop{:}}=\mathsf{ML}$. T-Elim is equivalent to Bot-E, so $\mathsf{ML}+\mathsf{T}$ corresponds to $\mathsf{IL}$. $\hfill\blacktriangleleft$\par
 
% END S5.SS0.SSS0.P0.SPx2.p1
 
% END S5.SS0.SSS0.P0.SPx2
 
% END S5


% BEGIN S6
 
\par\medskip\noindent\textbf{6 Fragment of Coherent Formulas}\par
 
% BEGIN S6.p1
 
Coherent formulas are formulas that can only be built using conjunctions, disjunctions and existential quantifiers. Because monad operators are only inserted after universal quantifiers and implications, it follows that $A_{T}=A$ for any coherent formula $A$. We lift to higher-order logic two tricks for coherent formulas that rely on this observation.\par
 
% END S6.p1
 
% BEGIN S6.p2
 
The first trick [5, 22] allows us to derive minimal provability from intuitionistic provability when the context is empty and when the formula is coherent.\par
 
% END S6.p2
 
% BEGIN Thmtheorem14
 
\par\medskip\noindent\textbf{Corollary 14.}\par
 
% BEGIN Thmtheorem14.p1
 
Let $\Gamma,A$ be higher-order coherent formulas. If $\vdash_{\mathsf{IL}}A$ then $\vdash_{\mathsf{ML}}A$.\par
 
% END Thmtheorem14.p1
 
% END Thmtheorem14
 
% BEGIN S6.SS0.SSS0.P0.SPx1
 
\par\medskip\noindent\textbf{Proof.}\par
 
% BEGIN S6.SS0.SSS0.P0.SPx1.p1
 
Suppose $\vdash_{\mathsf{IL}}A$. We apply Theorem 8 with $TA\mathrel{\mathop{:}}=A\mathbin{\vee}\bot$, and we derive $\vdash_{\mathsf{ML}}TA_{T}$, that is $\vdash_{\mathsf{ML}}A\mathbin{\vee}\bot$. We know that minimal logic has the disjunction property, and that we cannot prove $\bot$. Hence, we conclude $\vdash_{\mathsf{ML}}A$. $\hfill\blacktriangleleft$\par
 
% END S6.SS0.SSS0.P0.SPx1.p1
 
% BEGIN S6.SS0.SSS0.P0.SPx1.p2
 
The second trick [18, 22] allows us to derive intuitionistic provability from classical provability when the formulas involved are coherent. It corresponds to a weaker version–with finite disjunctions–of Barr’s theorem [2] extended to higher-order logic.\par
 
% END S6.SS0.SSS0.P0.SPx1.p2
 
% BEGIN Thmtheorem15
 
\par\medskip\noindent\textbf{Theorem 15.}\par
 
% BEGIN Thmtheorem15.p1
 
Let $\Gamma,A$ be higher-order coherent formulas. If $\Gamma\vdash_{\mathsf{CL}}A$ then $\Gamma\vdash_{\mathsf{IL}}A$.\par
 
% END Thmtheorem15.p1
 
% END Thmtheorem15
 
% BEGIN S6.SS0.SSS0.P0.SPx1.p3
 
To prove this result, we adapt to higher-order logic Palmgren’s idea [18] of using Friedman’s translation [9]. As an intermediary lemma, we prove that, for higher-order logic, Friedman’s translation transforms classical proofs into intuitionistic proofs. Remark that we cannot directly apply Theorem 8, as the $T$-elimination principle does not correspond to a classical law when $TA\mathrel{\mathop{:}}=(A\mathbin{\Rightarrow}R)\mathbin{\Rightarrow}R$.\par
 
% END S6.SS0.SSS0.P0.SPx1.p3
 
% BEGIN Thmtheorem16
 
\par\medskip\noindent\textbf{Lemma 16.}\par
 
% BEGIN Thmtheorem16.p1
 
Let $TA\mathrel{\mathop{:}}=(A\mathbin{\Rightarrow}R)\mathbin{\Rightarrow}R$ with parameter $R$. Let $\Gamma,A$ be higher-order formulas. If $\Gamma\vdash_{\mathsf{CL}}A$ then $\Gamma_{T}\vdash_{\mathsf{IL}}A^{T}$.\par
 
% END Thmtheorem16.p1
 
% END Thmtheorem16
 
% END S6.SS0.SSS0.P0.SPx1
 
% BEGIN S6.SS0.SSS0.P0.SPx2
 
\par\medskip\noindent\textbf{Proof.}\par
 
% BEGIN S6.SS0.SSS0.P0.SPx2.p1
 
We proceed by induction on the derivation. All the cases are similar to the cases of Theorem 8, except that we do not need the T-Elim case anymore, and that we change the PEM case. We prove $\Gamma_{T}\vdash_{\mathsf{IL}}T(A\mathbin{\vee}\neg A)_{T}$ using\par
 
% END S6.SS0.SSS0.P0.SPx2.p1
 
% BEGIN S6.SS0.SSS0.P0.SPx2.p2
 
\begin{center}\includegraphics[width=0.95\linewidth,height=0.45\textheight,keepaspectratio]{prior-graphic-04.pdf}\end{center}
 
% END S6.SS0.SSS0.P0.SPx2.p2
 
% BEGIN S6.SS0.SSS0.P0.SPx2.p3
 
where $B$ is an abbreviation for the formula $A_{T}\mathbin{\vee}(A_{T}\mathbin{\Rightarrow}((\bot\mathbin{\Rightarrow}R)\mathbin{\Rightarrow}R))$. $\hfill\blacktriangleleft$\par
 
% END S6.SS0.SSS0.P0.SPx2.p3
 
% END S6.SS0.SSS0.P0.SPx2
 
% BEGIN S6.SS0.SSS0.P0.SPx3
 
\par\medskip\noindent\textbf{Proof of Theorem 15.}\par
 
% BEGIN S6.SS0.SSS0.P0.SPx3.p1
 
Suppose $\Gamma\vdash_{\mathsf{CL}}A$. Without loss of generality, we assume that $\Gamma$ and $A$ have no free variables, as we can replace them by fresh constants. We apply Lemma 16, and we derive $\Gamma_{T_{R}}\vdash_{\mathsf{IL}}T_{R}A_{T_{R}}$, that is $\Gamma\vdash_{\mathsf{IL}}(A\mathbin{\Rightarrow}R)\mathbin{\Rightarrow}R$. Choosing $R\mathrel{\mathop{:}}=A$, we get $\Gamma\vdash_{\mathsf{IL}}A$. $\hfill\blacktriangleleft$\par
 
% END S6.SS0.SSS0.P0.SPx3.p1
 
% BEGIN S6.SS0.SSS0.P0.SPx3.p2
 
In other words, coherent formulas correspond to a constructive fragment of higher-order classical logic. Constructive fragments of classical logic have been studied for propositional logic [11, 12] and first-order logic [19]. This is the first time, to our knowledge, that such results are proved for higher-order logic.\par
 
% END S6.SS0.SSS0.P0.SPx3.p2
 
% END S6.SS0.SSS0.P0.SPx3
 
% END S6
